
\documentclass[11pt]{article}

\usepackage[margin=1in]{geometry}
\usepackage{amsmath,amssymb,amsthm}
\usepackage{array}
\usepackage[T1]{fontenc}
\usepackage{lmodern}
\usepackage[hidelinks]{hyperref}

\newtheorem{theorem}{Theorem}
\newtheorem{lemma}{Lemma}
\newtheorem{proposition}{Proposition}
\newtheorem{corollary}{Corollary}
\newtheorem{heuristic}{Heuristic}
\newtheorem{remark}{Remark}
\newtheorem{definition}{Definition}

\newcommand{\F}{\mathbb F}
\newcommand{\RS}{\operatorname{RS}}
\newcommand{\Bnum}{B_{\rm num}}
\newcommand{\Deltaone}{\Delta_1}
\newcommand{\Deltazero}{\Delta_0}
\newcommand{\tauone}{\tau_1}
\newcommand{\tautwo}{\tau_2}
\newcommand{\tauthree}{\tau_3}
\newcommand{\Ctwo}{C_2}
\newcommand{\qline}{q_{\rm line}}
\newcommand{\qgen}{q_{\rm gen}}
\newcommand{\qchal}{q_{\rm chal}}
\newcommand{\MCA}{\operatorname{MCA}}
\newcommand{\CA}{\operatorname{CA}}
\newcommand{\List}{\operatorname{List}}
\newcommand{\PG}{\operatorname{PG}}
\newcommand{\BCHKS}{\operatorname{BCHKS}}
\newcommand{\dist}{\Delta}
\newcommand{\efib}{c_{\rm fib}}
\newcommand{\ellpg}{\ell_{\rm pg}}

\title{Experimental Theorems and Imported-Locator Ledger for RS-MCA}
\author{RS-MCA experimental notes}
\date{June 18, 2026}

\begin{document}
\maketitle

\begin{abstract}
This note is the experimental collection point for statements that are useful
for the RS-MCA papers but are not yet main-paper theorems.  It has two jobs.
First, it records the correct provenance of the quotient-locator construction
used in a separate AI-generated packet: the locator proof underlying item (d)
is an imported theorem/proof pattern from Ben-Sasson--Carmon--Hab\"ock--
Kopparty--Saraf, \emph{On Proximity Gaps for Reed--Solomon Codes}
\cite{BCHKS25}, especially their Theorem~7.1 and the Theorem~1.13 subgroup
instantiation.  That proof is not a new contribution of this repository.
Second, it records genuinely local material that may be reusable later: the
smooth-quotient notation dictionary, a list-fiber pigeonhole wrapper around
the imported locator construction, the slack-two/subfield bookkeeping target
for Paper D, and a new divisibility gate for Paper B's quadratic-extension,
slack-two, three-co-support resonance window.

Everything here is experimental or imported unless explicitly promoted by a
main paper.  In particular, this note does not prove the corrected MCA
conjecture, does not give a protocol consequence, and does not supersede
Papers A--D.
\end{abstract}

\section{Status, Provenance, and What Counts as New}

This file is the place where we collect new things, but it must separate new
repository contributions from imported constructions.  The status labels used
below are:
\[
\begin{array}{ll}
\textsc{Imported} & \text{an external theorem, proof pattern, or conversion;}\\
\textsc{Wrapper} & \text{a repository repackaging, dictionary, or pigeonhole consequence;}\\
\textsc{New-local} & \text{a restricted algebraic theorem proved in this note;}\\
\textsc{Target} & \text{a precise proof obligation not yet promoted;}\\
\textsc{Heuristic} & \text{evidence or a search direction, not a theorem.}
\end{array}
\]

\paragraph{Hard attribution rule.}
Any locator-polynomial argument derived from the subgroup construction of
Ben-Sasson--Carmon--Hab\"ock--Kopparty--Saraf must cite \cite[Thm.~7.1]{BCHKS25}
where the construction is stated or proved.  Any admissible subgroup
instantiation must also cite \cite[Thm.~1.13]{BCHKS25}.  The citation belongs
at the theorem statement and again at the proof entry point; a bibliography
entry or acknowledgement alone is not enough.

\paragraph{What is new in this note.}
The following items are the repository-side contributions collected here.
\begin{itemize}
\item A smooth-quotient notation dictionary translating BCHKS subgroup
  notation into the RS-MCA notation \(Q=D^{\efib}\) and the dyadic/FFT-domain
  ledgers used in Papers A--D.
\item A list-fiber pigeonhole wrapper: once a locator construction assigns
  many subset certificates to a smaller slope set, one heavy slope has many
  locator certificates, and, after a distinctness check, many nearby
  codewords.
\item A slack-two and subfield-pigeonhole proof target for the Paper D route:
  the imported locator construction should be run at the augmented-code rung
  needed by the Crites--Stewart list-to-agreement conversion, and slopes should
  be pigeonholed over the generated/base field when the construction really
  takes values there.
\item A \textsc{New-local} divisibility gate in the restricted Paper B window
  \(B=\F_p\), \(F=\F_{p^2}\), \(D=\F_p\), \(t=\sigma=2\), \(j=3\).  This is
  independent of the BCHKS priority issue and is the main theorem proved in
  Section~\ref{sec:divisibility-gate}.
\end{itemize}

\paragraph{What is not new.}
The polynomial identity behind the item-(d) locator proof is not new here.  In
BCHKS notation it is the expansion of a vanishing polynomial over a complement
of a chosen subset; in smooth-quotient notation it is the same expansion of
\(\prod_{b\in A}(X^{\efib}-b)\), up to choosing a subset or its complement.
The use of \(\ell\) for the chosen subset size also follows the external
presentation closely enough that it should either be cited or renamed to
\(\ellpg\).

\paragraph{Source map.}
The main papers remain the authoritative documents.
\begin{itemize}
\item Paper A, \texttt{tex/RS\_disproof\_v3.tex}, refutes the old no-slack
  smooth-domain support-wise line-MCA statement.  If it uses the locator
  construction, it should cite BCHKS and call the repository contribution a
  smooth-domain/list-fiber corollary, not a new locator proof.
\item Paper B, \texttt{tex/slackMCA\_v3.tex}, is the main target for the local
  resonance material in Sections~\ref{sec:resonance-setup}--\ref{sec:scanner}.
\item Paper C, \texttt{tex/snarks\_v4.tex}, records protocol ledgers.  Nothing
  in this note may be consumed by a protocol without a separate conversion
  theorem.
\item Paper D, \texttt{tex/cs25\_cap\_v4.tex}, gives the universal cap via an
  imported list-to-agreement route.  Its locator-side input should cite BCHKS
  explicitly if it uses the quotient-locator construction.
\end{itemize}

\section{BCHKS Locator Construction in Smooth-Quotient Notation}
\label{sec:bchks}

\textbf{Status: \textsc{Imported} plus \textsc{Wrapper}.}
This section is a provenance-safe restatement and dictionary.  It is included
so later main-paper edits have a clean place to copy from.

\subsection{Imported theorem/proof pattern}

BCHKS prove a general multiplicative-subgroup construction
\cite[Thm.~7.1]{BCHKS25}.  In a compact form, the imported statement is as
follows.

\begin{theorem}[BCHKS quotient-locator construction, imported]
\label{thm:bchks-import}
Let \(G\subseteq H\subseteq \F_q^*\) be multiplicative subgroups and let
\(\Phi:H\to G\) be \(x\mapsto x^c\), where \(c=|H|/|G|\).  Let
\(E\subseteq G\), and let \(\ellpg,a\geq1\) satisfy
\[
  |E^{(+\ellpg)}|
  =\#\{e_1+\cdots+e_{\ellpg}:e_i\in E\text{ distinct}\}\geq a.
\]
Let \(D=\Phi^{-1}(E)\), \(n=|D|\), and
\[
  C=\RS[\F_q,D,n-(\ellpg+2)c].
\]
Then, for \(\gamma=\ellpg c/n\), there are functions
\(f,g:D\to\F_q\) such that
\[
  \#\{z\in\F_q: \dist(f+zg,C)\leq \gamma\}\geq a,
\]
while
\[
  \dist([f,g],C^2)\geq (\ellpg+1)c/n.
\]
\end{theorem}

\begin{proof}[Reference]
This is exactly the imported construction of \cite[Thm.~7.1]{BCHKS25}; the
proof is not reproduced as a repository contribution.  BCHKS define the line
endpoints by the monomials \(x^{n-\ellpg c}\) and
\(x^{n-(\ellpg+1)c}\), then expand a vanishing polynomial over the complement
of a chosen \(\ellpg\)-subset.  The coefficient of
\(X^{n-(\ellpg+1)c}\) is an affine transform of a restricted sum from
\(E^{(+\ellpg)}\), giving the exceptional slope values.
\end{proof}

The BCHKS Theorem~1.13 is the admissible-subgroup specialization of this
construction: take \(E=G\), take \(H=D\), set \(c=|D|/|G|\), and use
\(\ellpg=|G|/2\).  Any RS-MCA statement corresponding to that specialization
should cite \cite[Thm.~1.13]{BCHKS25} directly.

\subsection{Notation dictionary}

\begin{center}
\begin{tabular}{p{0.28\linewidth}p{0.30\linewidth}p{0.32\linewidth}}
\textbf{BCHKS notation} & \textbf{RS-MCA notation} & \textbf{Meaning}\\
\hline
\(H\) & evaluation domain \(D\) & Large multiplicative domain on which codewords are evaluated.\\
\(G\) & quotient \(Q=D^{\efib}\) & Image of the power map; one point in \(Q\) has \(\efib\) preimages in \(D\).\\
\(c=|H|/|G|\) & \(\efib=n/|Q|\) & Fiber size of the quotient map.  Avoid calling this \(a\), because BCHKS also uses \(a\) for the number of exceptional slopes.\\
\(\Phi(x)=x^c\) & \(x\mapsto x^{\efib}\) & Quotient/fiber map.\\
\(E\subseteq G\) & active quotient subset \(E\subseteq Q\) & Quotient-level support used by the locator.\\
\(E^{(+\ellpg)}\) & restricted first-coefficient/slope set & Sums of distinct chosen quotient elements; slopes are affine transforms of these sums.\\
\(\ellpg\) & \(\ellpg\) or renamed subset size & Number of chosen quotient elements.  Use \(\ellpg\) in RS-MCA text to avoid collision with other \(\ell\)'s.\\
\(f=x^{n-\ellpg c}\), \(g=x^{n-(\ellpg+1)c}\) & locator line endpoints & Same monomial line after quotient notation is substituted.\\
\(P_S(X)=\prod_{\alpha\in S'}(X^c-\alpha)\) & \(\prod_{b\in A}(X^{\efib}-b)\), possibly after complementing \(A\) & The locator polynomial.\\
exceptional \(z\)'s & bad slopes / heavy list fibers & Slope parameters for which a line word is close to the RS code.\\
\end{tabular}
\end{center}

\subsection{The shared locator identity}

The following identity is written down because it is the exact place where
priority can be lost.  The identity is the BCHKS identity in RS-MCA notation.

\begin{lemma}[Locator expansion, imported identity]
\label{lem:locator-imported}
Let \(D\) map \(\efib\)-to-one onto a quotient set \(E\) by
\(x\mapsto x^{\efib}\), and let \(A\subseteq E\) have size \(\ellpg\).  Put
\(A^c=E\setminus A\) and define
\[
  P_A(X)=\prod_{\alpha\in A^c}(X^{\efib}-\alpha).
\]
Then
\[
  P_A(X)
  =X^{n-\ellpg\efib}
    -e_1(A^c)X^{n-(\ellpg+1)\efib}
    +\text{terms of degree at most }n-(\ellpg+2)\efib.
\]
Equivalently, since \(e_1(A^c)=e_1(E)-e_1(A)\), the coefficient of the next
monomial is an affine transform of the restricted sum \(e_1(A)\).
\end{lemma}

\begin{proof}[Reference]
This is the elementary expansion used inside \cite[Thm.~7.1]{BCHKS25}.  We
record it only to fix notation.  It should not be cited as a new lemma of the
repository unless the BCHKS provenance is present at the statement.
\end{proof}

\subsection{List-fiber pigeonhole wrapper}

The next proposition is the repository-side wrapper around the imported
identity.  It is intentionally stated in certificate form, because distinctness
of the nearby codewords is a separate check in concrete instantiations.

\begin{proposition}[List-fiber pigeonhole wrapper]
\label{prop:list-fiber-wrapper}
Let \(\mathcal A\) be a family of \(\ellpg\)-subsets of a quotient set \(E\).
Suppose the BCHKS locator construction assigns to each \(A\in\mathcal A\) a
slope \(z(A)\) lying in a finite set \(S\), and a locator certificate showing
that the corresponding word \(f+z(A)g\) is within radius \(\gamma\) of a
locator codeword \(P_A\in C\).  Then some slope \(z_0\in S\) has at least
\[
  \frac{|\mathcal A|}{|S|}
\]
locator certificates above it.  If the codewords \(P_A\) attached to the
certificates over \(z_0\) are distinct, then the Hamming ball of radius
\(\gamma\) around \(f+z_0g\) contains at least \(|\mathcal A|/|S|\) codewords
of \(C\).
\end{proposition}

\begin{proof}
The first assertion is the pigeonhole principle applied to the map
\(A\mapsto z(A)\).  The second assertion is only the definition of list size,
after the distinctness check removes possible duplicate locator codewords.
The construction that supplies the certificates is imported from
Theorem~\ref{thm:bchks-import}.
\end{proof}

\begin{remark}[What the wrapper adds]
The new part here is not the locator polynomial.  The wrapper isolates the
extra bookkeeping needed by RS-MCA: which slope set pays the pigeonhole bill,
whether the slope set is the ambient field \(\F\) or a generated/base subfield
\(B\), and whether locator certificates remain distinct after specialization.
Those are the checks needed before using the construction in Papers A, B, or D.
\end{remark}

\subsection{Slack-two and subfield target for Paper D}

\textbf{Status: \textsc{Target}.}
The Paper D route needs a locator-heavy list fiber in the augmented code rung
\(C^+=\RS[\F,D,k+1]\), because that is the object fed into the
Crites--Stewart list-to-agreement conversion.  The proof obligation should be
stated as follows.

\begin{proposition}[Slack-two/subfield locator target]
\label{prop:slack-two-target}
Let \(D\subseteq B\subseteq\F\), and suppose a BCHKS-type quotient locator
construction over \(D\) produces slopes in \(B\), not merely in \(\F\).  If the
one-rung/augmented-code version produces distinct locator codewords in
\(C^+=\RS[\F,D,k+1]\) for a family \(\mathcal A\) of certificates, then there
exists a slope \(z_0\in B\) with list size at least
\[
  \frac{|\mathcal A|}{|B|}
\]
inside the appropriate Hamming ball for \(C^+\).
\end{proposition}

\begin{proof}[Bookkeeping only]
This is Proposition~\ref{prop:list-fiber-wrapper} with \(S=B\).  The nontrivial
work is external or elsewhere: one must verify the BCHKS locator hypotheses at
the augmented-code rung, verify that the slope coefficients lie in \(B\), and
verify distinctness of the produced codewords.  Until those checks are present,
this proposition is a target, not a promoted Paper D theorem.
\end{proof}

\subsection{Pasteable main-paper wording}

\paragraph{Paper A wording.}
The polynomial locator identity used here is the quotient-fiber form of the
construction of Ben-Sasson--Carmon--Hab\"ock--Kopparty--Saraf
\cite[Thm.~7.1]{BCHKS25}; their Theorem~1.13 is the admissible subgroup
instantiation.  Thus the locator proof itself is not a new contribution of
this manuscript.  The present use is the smooth-domain specialization, the
field/divisor bookkeeping, and the list-fiber pigeonhole packaging.

\paragraph{Paper D wording.}
The locator-side input is the BCHKS quotient-locator construction
\cite[Thm.~7.1; cf.~Thm.~1.13]{BCHKS25}, written in the smooth-domain quotient
notation \(Q=D^{\efib}\).  The additional bookkeeping needed here is the
slack-two augmented-code rung and, when available, subfield pigeonholing over
\(B\) rather than over the ambient extension field \(\F\).

\section{Triage of the AI-Generated Four-Item Packet}
\label{sec:packet-triage}

\textbf{Status: \textsc{Audit}.}
The four advertised items should be positioned as follows before any promotion
into a main paper.  The original AI packet used the labels (a)--(d) without
making them self-contained in this note, so we decode them here.

\begin{center}
\begin{tabular}{p{0.10\linewidth}p{0.28\linewidth}p{0.50\linewidth}}
\textbf{Item} & \textbf{Short name} & \textbf{Meaning in this ledger}\\
\hline
(a) & Weak-slack positive regime & A positive-looking proximity/list or
MCA-adjacent statement that only reaches slack on the order of
\(1/\sqrt n\).  It is useful as a comparison theorem, but it is far from the
constant/reserve bookkeeping needed for the corrected RS-MCA target.\\
(b) & Finite Fermat-prime packet & A special finite-domain instantiation over
only three Fermat-prime cases, in the current experimental ledger represented
by the \(p\in\{17,257,65537\}\) style checks.  It is evidence or a worked
computation, not a uniform theorem.\\
(c) & Exponential-field construction & A large-field variant whose proof
requires the ambient field to be exponential in the domain/parameter size.
This can clarify the mechanism but does not match the polynomial-field or
challenge-field regimes needed by the Proximity Prize ledger.\\
(d) & BCHKS quotient-locator packet & The interesting locator/proximity-gap
construction: choose a quotient of a multiplicative subgroup, use
\(\ellpg\)-subsets and restricted sums to create many exceptional slopes or
nearby locator certificates.  This is the BCHKS construction of
\cite[Thm.~7.1]{BCHKS25}, with the admissible subgroup instance in
\cite[Thm.~1.13]{BCHKS25}; the repository contribution is only the
smooth-domain dictionary, list-fiber wrapper, and later slack-two/subfield
bookkeeping.\\
\end{tabular}
\end{center}

Item (d) should not be confused with the \textsc{New-local} divisibility-gate
theorem in Section~\ref{sec:divisibility-gate}.  The former is an imported
locator/proximity construction; the latter is an independent restricted Paper
B resonance calculation.

\begin{center}
\begin{tabular}{p{0.10\linewidth}p{0.34\linewidth}p{0.44\linewidth}}
\textbf{Item} & \textbf{Limitation} & \textbf{Repository handling}\\
\hline
(a) & Works only at slack on the order of \(1/\sqrt n\). & Use as a weak-slack comparison point only.  It does not clear the corrected-reserve MCA objective.\\
(b) & Applies only to three Fermat-prime cases. & Treat as finite-prime evidence or a worked computation, not as a uniform smooth-domain theorem.\\
(c) & Requires exponential field size. & Treat as a large-field sanity check.  It is not directly usable for polynomial-field or prize-ledger parameters.\\
(d) & Most relevant, but the locator proof is essentially BCHKS. & Import, cite, and adapt from \cite[Thm.~7.1 and Thm.~1.13]{BCHKS25}.  The repository contribution is only the smooth-domain/list-fiber/slack-two packaging after exact hypotheses are checked.\\
\end{tabular}
\end{center}

\paragraph{Promotion checklist.}
Before any statement using item (d) moves to Papers A--D, check the following.
\begin{itemize}
\item Does the statement cite BCHKS at the theorem and proof entry point?
\item Is the object a proximity-gap statement, a list statement, CA, MCA,
  line-decoding, or a protocol-facing theorem?
\item Which field pays the slope/list pigeonhole bill: \(B\), \(\F\), or
  another generated field?
\item Are the locator codewords distinct after specialization?
\item Is the claimed slack in the same normalization as the consumed theorem?
\item Are items (a), (b), and (c) kept out of asymptotic/protocol claims unless
  their limitations are explicitly part of the theorem statement?
\end{itemize}

\section{Restricted Resonance Setup for Paper B}
\label{sec:resonance-setup}

\textbf{Status: \textsc{New-local} for the divisibility gate; otherwise local
normal-form bookkeeping.}
The remaining sections preserve and reorganize the Cycle 14--18 algebraic
material.  This branch is independent of the BCHKS priority correction.

The common restricted setting is
\[
  B=\F_p,\qquad F=\F_{p^2},\qquad D=\F_p,\qquad
  t=\sigma=2,\qquad j=3,
\]
and we work off the tangent/global endpoint
\[
  R_0=\{\, [W]_E\wedge [\Bnum]_E=0\,\}.
\]
The field ledgers remain distinct:
\[
  \qgen=p,\qquad \qline=p^2.
\]
The regime is sub-reserve, since \(\eta=\sigma/n=2/n\).  Thus none of the
following proves a corrected-reserve list, CA, MCA, line-decoding, SNARK, or
protocol theorem.

The base/generated field is \(B=\F_p\).  The line field is the quadratic
extension \(F=\F_{p^2}=B(\alpha)\), with basis \(\{1,\alpha\}\) over \(B\) and
\(\alpha^2=\nu\in B\).  The domain is \(D=\F_p\), so \(n=|D|=p\).

A degree-two residue-line datum consists of
\[
  E\in F[X],\qquad \deg E=2,
  \qquad E(x)\ne0\ \text{for all }x\in D,
\]
and a numerator
\[
  \Bnum\in F[X],\qquad \deg \Bnum<2.
\]
We write
\[
  A=F[X]/E,
  \qquad [P]_E\in A,
  \qquad b=[\Bnum]_E.
\]
The anchor word on \(D\) is interpolated by a polynomial \(W\).  The operation
\(u\wedge v\) means the determinant of the coordinate vectors of \(u,v\in A\)
in a chosen \(F\)-basis of the two-dimensional \(F\)-space \(A\).  Off \(R_0\),
the pair \(\{[W]_E,b\}\) is such a basis.

The co-support has size \(j=3\).  If \(T\subset D\) is the three-point
co-support, define
\[
  \tau_1=e_1(T),\qquad \tau_2=e_2(T),\qquad \tau_3=e_3(T).
\]
All valid co-supports have \(\tau_i\in B\).  The quantity \(\Ctwo\) denotes the
number of distinct slopes produced by the relevant landing branch after
restricting to valid split triples \(T\subset\F_p\).  Bounding landing points
therefore bounds \(\Ctwo\), but a sharp slope theorem may require more because
many landing points can share a slope.

\section{Cycle 14 Normal Form}
\label{sec:cycle14}

The Cycle 14 audit reduces the \(t=2,j=3\) landing equation to two affine forms
in the elementary co-support parameters
\[
  \tau=(\tauone,\tautwo,\tauthree).
\]

\begin{lemma}[Affine wedge normal form]
\label{lem:affine-wedge}
Assume \(\{[W]_E,b\}\) is an \(F\)-basis of \(A\).  Then the landing equation can
be written in the form
\[
  \iota(\tau)=A_0(\tauone,\tautwo)-\tauthree [W]_E,
  \qquad
  \mu(\tau)=B_0(\tauone,\tautwo)-\tauthree b,
\]
where
\[
  A_0=p_1[W]_E+p_2 b,
  \qquad
  B_0=q_1[W]_E+q_2 b,
\]
and \(p_i,q_i\in F\) are affine-linear functions of
\((\tauone,\tautwo)\).  If the basis is normalized so \([W]_E\wedge b=1\), then
\[
  \Delta=\iota\wedge\mu
    =(p_1-\tauthree)(q_2-\tauthree)-p_2q_1.
\]
\end{lemma}

\begin{proof}
The Cycle 14 quotient calculation, which is a concrete instance of Paper B's
residue-line normal form, gives the displayed affine forms.  Expanding in the
ordered basis \(\{[W]_E,b\}\), the coordinate vector of \(\iota\) is
\((p_1-\tauthree,p_2)\), while the coordinate vector of \(\mu\) is
\((q_1,q_2-\tauthree)\).  Their wedge determinant is therefore
\[
  (p_1-\tauthree)(q_2-\tauthree)-p_2q_1.
\]
\end{proof}

\section{Cycle 18 Component Split}
\label{sec:cycle18}

Choose a basis \(F=B+\alpha B\) with \(\alpha^2=\nu\in B\), and write
\[
  p_1=p_{10}+\alpha p_{11},\quad
  p_2=p_{20}+\alpha p_{21},\quad
  q_1=q_{10}+\alpha q_{11},\quad
  q_2=q_{20}+\alpha q_{21}.
\]
Split
\[
  \Delta=\Deltazero+\alpha\Deltaone,
  \qquad
  \Deltazero,\Deltaone\in B[\tauone,\tautwo,\tauthree].
\]

\begin{theorem}[Monic quadratic and linear alpha component]
\label{thm:component-split}
In the setting above,
\[
\begin{aligned}
\Deltazero={}&
  \tauthree^2-(p_{10}+q_{20})\tauthree
  +p_{10}q_{20}+\nu p_{11}q_{21}
  -p_{20}q_{10}-\nu p_{21}q_{11},\\
\Deltaone={}&
  -(p_{11}+q_{21})\tauthree
  +p_{10}q_{21}+p_{11}q_{20}-p_{20}q_{11}-p_{21}q_{10}.
\end{aligned}
\]
Consequently \(\Deltazero\) is monic quadratic in \(\tauthree\), while
\(\Deltaone\) has degree at most one in \(\tauthree\).
\end{theorem}

\begin{proof}
Substitute the coordinate expressions for \(p_i,q_i\) into the identity of
Lemma~\ref{lem:affine-wedge} and use \(\alpha^2=\nu\).  Collecting the
coefficient of \(1\) gives the displayed formula for \(\Deltazero\), and
collecting the coefficient of \(\alpha\) gives the displayed formula for
\(\Deltaone\).  The coefficient of \(\tauthree^2\) in \(\Deltazero\) is \(1\),
and the \(\alpha\)-component has no \(\tauthree^2\)-term.
\end{proof}

\begin{corollary}[Graph reduction outside the zero-alpha branch]
\label{cor:graph}
If \(\Deltaone\) is not the zero polynomial, write
\[
  \Deltaone=s(\tauone,\tautwo)\tauthree+h(\tauone,\tautwo).
\]
On the open locus \(s\ne0\), the common-zero condition
\(\Deltazero=\Deltaone=0\) is supported on the graph
\[
  \tauthree=-h/s.
\]
\end{corollary}

\begin{proof}
This is immediate from solving the linear equation \(\Deltaone=0\) for
\(\tauthree\) on \(s\ne0\).
\end{proof}

\section{New Divisibility-Gate Theorem}
\label{sec:divisibility-gate}

The next theorem is the main \textsc{New-local} theorem extracted from the
Cycle 18 reduction.  It is stronger than merely saying that the branch is a
graph: it shows that any two-dimensional family must satisfy an explicit
divisibility identity.

\begin{theorem}[Divisibility gate for the graph branch]
\label{thm:divisibility-gate}
Assume the restricted setting above and assume \(\Deltaone\ne0\).  Write
\[
  \Deltaone=s(\tauone,\tautwo)\tauthree+h(\tauone,\tautwo),
\]
and define the denominator-cleared graph polynomial
\[
  G(\tauone,\tautwo)
  =s(\tauone,\tautwo)^2
    \Deltazero\bigl(\tauone,\tautwo,-h(\tauone,\tautwo)/s(\tauone,\tautwo)\bigr).
\]
Then \(\deg G\le 4\).  If \(G\) is not the zero polynomial, then
\[
  \#\{\,\tau\in B^3:\Deltazero(\tau)=\Deltaone(\tau)=0\,\}=O(p).
\]
Consequently the corresponding number of distinct slopes satisfies
\[
  \Ctwo=O(p)=O(n).
\]
Therefore a growing-prime family with
\[
  \Ctwo=\Theta(p^2)=\Theta(\qline)
\]
in the nonzero-\(\Deltaone\) branch can occur only if \(G\equiv0\).
\end{theorem}

\begin{proof}
Because the \(p_i,q_i\) are affine-linear in \((\tauone,\tautwo)\), the
coefficient \(s\) has degree at most \(1\), and \(h\) has degree at most \(2\).
Since \(\Deltazero\) is monic quadratic in \(\tauthree\), clearing the
denominators after substituting \(\tauthree=-h/s\) gives a polynomial \(G\) of
degree at most \(4\).

On \(s\ne0\), the equation \(\Deltaone=0\) forces \(\tauthree=-h/s\).  Hence
every common zero on this open locus gives a base pair
\((\tauone,\tautwo)\) with \(G(\tauone,\tautwo)=0\).  If \(G\) is nonzero,
Schwartz--Zippel gives at most \(4p\) such base pairs, and each base pair gives
at most one value of \(\tauthree\).

It remains to count the exceptional locus \(s=0\).  There \(\Deltaone=0\) also
forces \(h=0\).  If \(s\) is a nonzero affine form, this is contained in an
affine line; if \(s\equiv0\), then \(h\) is a nonzero polynomial of degree at
most \(2\).  In either case there are \(O(p)\) base pairs.  For each such pair,
\(\Deltazero\) is a monic quadratic in \(\tauthree\), so it gives at most two
values of \(\tauthree\).  The total common-zero count is \(O(p)\).

Finally, the number of distinct slopes is bounded by the number of landing
points \(\tau\).  Thus \(\Ctwo=O(p)\) whenever \(G\not\equiv0\), and any
\(\Theta(p^2)\) family in this branch must satisfy \(G\equiv0\).
\end{proof}

\begin{remark}
The identity \(G\equiv0\) says that, after clearing the \(s\)-denominator,
\(\Deltazero\) is divisible by the linear equation \(\Deltaone\) on the graph
branch.  This is the first exact algebraic gate that a large slope family in
this restricted window must pass.
\end{remark}

\section{Nearby Proven Gates}
\label{sec:nearby-gates}

The same restricted lane contains two other useful gates.  They are included
here because they explain where Theorem~\ref{thm:divisibility-gate} fits.

\begin{proposition}[Base-component complete-intersection gate]
\label{prop:complete-intersection}
Suppose \(\Deltazero,\Deltaone\in B[\tauone,\tautwo,\tauthree]\) have no common
surface component.  Then
\[
  \#V(\Deltazero,\Deltaone)(B)=O(p),
\]
and hence this branch contributes \(\Ctwo=O(p)\) slopes.
\end{proposition}

\begin{proof}
The two components have bounded degree.  If they share no surface component,
their common zero set in \(\mathbb A^3_B\) is a curve of bounded degree.  A
bounded-degree affine curve over \(\F_p\) has \(O(p)\) rational points, and
slopes are bounded by landing points.
\end{proof}

\begin{proposition}[Cycle 16 determinant gate]
\label{prop:det-gate}
Let \(Q(z_0,z_1)\) be the Cycle 15--16 determinant consistency polynomial for
the affine system
\[
  L_z(\tau)=\iota(\tau)-z\mu(\tau)=0,
  \qquad z=z_0+\alpha z_1.
\]
If \(Q\) is not identically zero, then in the restricted \(D=\F_p\),
\(t=\sigma=2\), \(j=3\), off-\(R_0\) window,
\[
  \Ctwo\le 4p=O(p)=O(n).
\]
\end{proposition}

\begin{proof}
Cycle 16 gives \(\deg Q\le4\) and shows that every landing slope must satisfy
\(Q(z_0,z_1)=0\).  If \(Q\ne0\), Schwartz--Zippel over \(B=\F_p\) gives at most
\(4p\) possible pairs \((z_0,z_1)\), hence at most \(4p\) slopes.
\end{proof}

\section{Slope Map and Remaining Heuristics}
\label{sec:slope-map}

On the graph branch, the landing equation \(\iota=z\mu\) gives
\[
  p_1-\tauthree=zq_1,
  \qquad
  p_2=z(q_2-\tauthree).
\]
Where \(q_1\ne0\), substituting \(\tauthree=-h/s\) yields
\[
  z(\tauone,\tautwo)=\frac{p_1+h/s}{q_1}.
\]
Equivalently, eliminating \(\tauthree\) recovers the Cycle 14 quadratic
\[
  q_1z^2-(p_1-q_2)z-p_2=0,
\]
so the residual map is controlled by
\[
  (\tauone,\tautwo)
  \longmapsto [\,q_1:(p_1-q_2):p_2\,]\in\mathbb P^2(F).
\]

\begin{heuristic}[Exceptional identity should be rare]
For generic source data, the polynomial \(G\) should not vanish identically.
Thus the graph branch should usually be curve-sized, and the first possible
large-slope families should come from algebraically special data satisfying
the divisibility identity \(G\equiv0\) or from the zero-alpha branch
\(\Deltaone\equiv0\).
\end{heuristic}

\begin{heuristic}[What remains after the gate]
If \(G\equiv0\), the problem is no longer to show that the common-zero set is
small.  It may be surface-sized.  The remaining question is whether the
projective coefficient map has one-dimensional image on every source-valid
surface, giving \(\Ctwo=O(p)\), or whether some growing-prime source-valid
family has two-dimensional image and \(\Ctwo=\Theta(p^2)\).
\end{heuristic}

\section{Next Experimental Checks}
\label{sec:scanner}

A useful scanner should compute, for each source-valid candidate,
\[
  \Deltaone=s\tauthree+h,
  \qquad
  G=s^2\Deltazero(\tauone,\tautwo,-h/s),
\]
then record whether \(G\equiv0\), whether \(\Deltaone\equiv0\), the exceptional
locus \(s=h=0\), the image size of
\[
  [\,q_1:(p_1-q_2):p_2\,],
\]
and the direct slope count over distinct split triples \(T\subset\F_p\).

The certificate should keep at least
\[
\begin{gathered}
  p,\ \qgen,\ \qline,\ E,\ \Bnum,\ \text{seed},\ \text{stratum},\
  R_0\text{-status},\\
  \Deltaone\equiv0,
  \quad G\equiv0,
  \quad \#\{s=h=0\},
  \quad \#\operatorname{image},
  \quad \Ctwo,
  \quad \#\{T\text{ split distinct}\}.
\end{gathered}
\]

\section{Non-Claims}
\label{sec:nonclaims}

This note does not claim:
\begin{itemize}
\item a proof of the corrected MCA conjecture;
\item a positive theorem above corrected reserve;
\item a \(\qgen\) local-limit theorem;
\item a list, CA, MCA, line-decoding, SNARK, or protocol consequence;
\item a large \(\Theta(p^2)\) counterpacket;
\item priority for the BCHKS quotient-locator proof pattern.
\end{itemize}

The mathematical value is narrower.  On the imported side, the file gives a
provenance-safe dictionary and wrapper for using BCHKS without claiming it as
new.  On the local Paper B side, the possible large branch in the restricted
quadratic-extension window must satisfy an explicit algebraic identity before
it can even be surface-sized.

\appendix

\section{BibTeX Entry to Copy into Main Papers}

\begin{verbatim}
@misc{BCHKS25,
  author = {Eli Ben-Sasson and Dan Carmon and Ulrich Habock and
            Swastik Kopparty and Shubhangi Saraf},
  title  = {On Proximity Gaps for Reed--Solomon Codes},
  year   = {2025},
  note   = {Manuscript dated November 11, 2025},
  url    = {https://www.math.toronto.edu/swastik/rs-proximity-gaps-2025.pdf}
}
\end{verbatim}

\begin{thebibliography}{9}
\bibitem{BCHKS25}
Eli Ben-Sasson, Dan Carmon, Ulrich Hab\"ock, Swastik Kopparty, and Shubhangi
Saraf.  \emph{On Proximity Gaps for Reed--Solomon Codes}.  Manuscript dated
November 11, 2025.
\url{https://www.math.toronto.edu/swastik/rs-proximity-gaps-2025.pdf}.
\end{thebibliography}

\end{document}
